\section{Case Study}\label{sec:case}

A case study on a small dataset checks the bench and the oracle and shows how the strategies differ on concrete
conflicts. All tables and the figure are generated from the simulator output by a script.

\subsection{Dataset and scenarios}
The dataset contains the six demonstration subscriptions of the application (streaming, music, gym, phone plan,
language course and cloud storage), with fixed UUIDs and dates, one or two tags each, one payment each and a
shared-with counter. Three clients $c_1, c_2, c_3$ start from the same server state. Six scripted scenarios each
isolate one kind of conflict:
C1: offline, $c_1$ edits the notes and $c_2$ the category of one record; both edits should survive.
C2: $c_1$ and $c_2$ both change the price offline (16000 and 14500 cents), a true conflict that should reach the user.
C3: $c_1$ deletes a record that $c_2$ edits offline, a policy choice. C4: $c_2$'s clock is one hour behind; $c_1$
renames a record, $c_2$ receives the rename and renames it again, so $c_2$'s causally later name should win. C5: the
three clients each add a different tag offline and $c_1$ also removes one; the result should be the union of the
changes. C6: the three clients each increment the shared-with counter offline, from 1 to an expected 4.

\subsection{Results of the scripted scenarios}
Table~\ref{tab:scenarios} gives the outcome of the scenario operations for each strategy, as judged by the oracle
(Definition~\ref{def:outcome}), and the total payload of the run. All strategies converged in all scenarios.

\begin{table}[ht]
\caption{Scripted scenarios C1--C6: lost and surfaced operations, policy applied to the ambiguous pair of C3, and
total payload. An asterisk marks a causal violation (a stale value won).}
\label{tab:scenarios}
\footnotesize
\setlength{\tabcolsep}{4pt}
% Generated by `npm run tables` (sim/scripts/tables.ts) from sim/out/case-study; do not edit.

\begin{tabular}{@{}lrrrrr@{}}
\toprule
 & LWW & VV & Automerge & Yjs & Hybrid \\
\midrule
\multicolumn{6}{@{}l}{\emph{Outcome of the scenario operations}} \\
C1 & 1 lost$^{\ast}$ & none & none & none & none \\
C2 & 1 lost & 1 surf. & 1 lost & 1 lost & 1 surf. \\
C3 & edit wins & conflict & delete wins & delete wins & edit wins \\
C4 & 1 lost$^{\ast}$ & none & none & none & none \\
C5 & 3 lost & 3 surf. & none & none & none \\
C6 & 2 lost & 2 lost & none & none & none \\
\midrule
\multicolumn{6}{@{}l}{\emph{Payload (bytes)}} \\
C1 & 2073 & 5495 & 1744 & 387 & 1147 \\
C2 & 1988 & 5542 & 1670 & 298 & 1259 \\
C3 & 1438 & 5051 & 1777 & 222 & 1017 \\
C4 & 2621 & 5414 & 1507 & 370 & 1076 \\
C5 & 2521 & 8234 & 2825 & 480 & 1318 \\
C6 & 2600 & 8185 & 2637 & 528 & 1446 \\
\bottomrule
\end{tabular}

\end{table}

LWW loses an edit in every scenario except C3. In C1 the loss needs no real conflict: whole-record replacement drops
$c_1$'s notes, and the visible notes is the old value $c_1$ had replaced (a causal violation). C4 shows the effect of clock skew that RQ1 asks about: $c_2$'s
causally later rename carries a timestamp one hour too early, so the stale name wins. Version vectors (S2) lose
nothing in C1--C5, but they treat tags and the counter as plain values: C5 produces a conflict with three lists the
user has to merge by hand, and in C6 the three candidates are equal (each client computed $1+1$), so two increments
disappear without a conflict. Automerge and Yjs merge C1, C4, C5 and C6 ideally but choose a winner silently in C2
(each library its own: Automerge keeps 14500, Yjs 16000) and let the delete win in C3. Only the hybrid strategy reaches
the ideal outcome in all six scenarios: C2 is surfaced, C3 keeps the edited record, the HLC protects C4 from skew, and
the set and counter merge exactly. Its payload is about half that of LWW, whose messages repeat whole
records.

\subsection{Small random workload}
The second part of the case study is a randomised run: three clients, six records, two simulated hours, about ten
operations per client per hour with all operation kinds, online and offline periods of 20 and 15 minutes on average,
a synchronisation every minute while online, and a skew of $-60$~s on $c_3$; five seeds.
Table~\ref{tab:random} and Figure~\ref{fig:random} give the results and Table~\ref{tab:byclass} breaks the losses down
by field class.

\begin{table}[ht]
\caption{Small random workload: mean $\pm$ sample standard deviation over five seeds (about 63 operations per run).
1~KB = 1000 bytes; the state is that of the server replica.}
\label{tab:random}
\footnotesize
\setlength{\tabcolsep}{3pt}
% Generated by `npm run tables` (sim/scripts/tables.ts) from sim/out/case-study; do not edit.

\begin{tabular}{@{}lrrrrrr@{}}
\toprule
 & Lost (\%) & Surf.\ (\%) & Viol. & Payload (KB) & State (KB) & Conv. \\
\midrule
LWW & $17.7 \pm 9.3$ & $0.0 \pm 0.0$ & $5.8 \pm 3.2$ & $54.6 \pm 18.5$ & $3.8 \pm 0.5$ & 5/5 \\
VV & $1.0 \pm 0.9$ & $3.2 \pm 2.9$ & $0.0 \pm 0.0$ & $146.9 \pm 33.1$ & $9.2 \pm 0.8$ & 5/5 \\
Automerge & $1.2 \pm 1.2$ & $0.0 \pm 0.0$ & $0.0 \pm 0.0$ & $44.3 \pm 11.5$ & $3.0 \pm 0.4$ & 5/5 \\
Yjs & $1.2 \pm 1.2$ & $0.0 \pm 0.0$ & $0.0 \pm 0.0$ & $21.0 \pm 7.7$ & $3.7 \pm 0.6$ & 5/5 \\
Hybrid & $0.6 \pm 0.9$ & $0.6 \pm 0.9$ & $0.0 \pm 0.0$ & $32.3 \pm 7.3$ & $6.2 \pm 0.6$ & 5/5 \\
\bottomrule
\end{tabular}

\end{table}

\begin{table}[ht]
\caption{Lost operations per field class over the five seeds of the random workload (lost / judged).}
\label{tab:byclass}
\footnotesize
% Generated by `npm run tables` (sim/scripts/tables.ts) from sim/out/case-study; do not edit.

\begin{tabular}{@{}lrrrrrr@{}}
\toprule
 & SIMPLE & CRITICAL & SET & APPEND & COUNTER & RECORD \\
\midrule
LWW & 20/124 & 10/46 & 14/53 & 7/35 & 3/31 & 0/25 \\
VV & 0/124 & 0/46 & 2/53 & 1/35 & 0/31 & 0/25 \\
Automerge & 2/124 & 2/46 & 0/53 & 0/35 & 0/31 & 0/25 \\
Yjs & 2/124 & 2/46 & 0/53 & 0/35 & 0/31 & 0/25 \\
Hybrid & 2/124 & 0/46 & 0/53 & 0/35 & 0/31 & 0/25 \\
\bottomrule
\end{tabular}

\end{table}

\begin{figure}[ht]
\centering
\begin{tikzpicture}
\begin{axis}[
  width=0.85\textwidth, height=4.6cm,
  xlabel={total payload (KB)}, ylabel={lost operations (\%)},
  xmin=0, xmax=170, ymin=-1, ymax=30,
  every axis label/.append style={font=\footnotesize},
  tick label style={font=\footnotesize}]
\addplot[only marks, mark=*, mark options={fill=black!40, draw=black},
  nodes near coords, point meta=explicit symbolic,
  every node near coord/.append style={font=\scriptsize, rotate=55, anchor=west, xshift=1pt}]
  table[x=payloadKB, y=lossPct, meta=label] {generated/random.dat};
\end{axis}
\end{tikzpicture}
\caption{Small random workload, mean over five seeds: share of lost operations against total payload. The ideal
strategy is in the lower left corner.}
\label{fig:random}
\end{figure}

LWW loses about one operation in six, in every field
class except record creation and deletion, and produces causal violations because of the skewed client. S2 loses almost nothing but surfaces the most
conflicts and sends the most bytes, because every message carries version vectors and dots for whole records.
Automerge and Yjs lose only the concurrent writes of scalar fields that their registers resolve silently, including
critical fields. S4 has the lowest loss: its two lost operations over all seeds are concurrent writes of SIMPLE fields,
lost by design to HLC-ordered LWW, while every concurrent change of a critical field is surfaced. Its payload is 41\%
smaller than that of LWW and 27\% smaller than that of Automerge, but larger than that of Yjs, whose binary encoding is
the most compact. The cost of S4 is its state: every field keeps its stamp and every record its causal context and
tombstone, so its server state is above that of LWW and Automerge. Compacting contexts that all replicas have
acknowledged is the first optimisation to try.

Five runs of one small configuration support no statistical claim, but the expectations of
Table~\ref{tab:expectations} already hold for loss and payload, not yet for state size.
